\section{Financial Literacy Through the Learn Module}

The Learn module is the platform's primary contribution to SDG 4 and serves as the entry point to the integrated learning pipeline. It delivers a structured curriculum of 23 financial concepts organized into five branches: Money Basics (income, expenses, budgeting), Saving (saving money, emergency funds, inflation), Growing Money (simple interest, compound interest, risk and return), Investing (bank deposits, fixed and recurring deposits, stocks, mutual funds, index funds, SIPs, diversification, fees and expense ratios, investment time horizon, risk tolerance) and Financial Safety (UPI safety, digital financial fraud, borrowing and interest, credit scores).

Each concept is served from a versioned lesson file that contains objectives, plain-language explanations, worked examples, key takeaways, a three-question quiz with explanations per objective, citations to verified sources and a review record. The curriculum prerequisite graph ensures that a concept unlocks only when all prerequisites are completed, and locked concepts display the missing prerequisite. Completion requires correct answers to all quiz questions.

The curriculum covers financial terminology and concepts, savings and budgeting, inflation and compound interest, savings accounts and fixed deposits, mutual funds, SIPs and stocks, investment risk and financial awareness. All content is grounded in verified publications: the RBI FAME fourth edition, SEBI investor education documents including Understanding Mutual Funds, Index Mutual Funds, How to Manage Investment Risks, Understanding Shares and Key Risks in Investing in Securities Market, AMFI disclosures on SIPs and Expense Ratio, and the DICGC information leaflet. Each lesson records its sources, access dates, content version and review status, with newly authored lessons labelled as source-checked pending human review.

Interactive learning and quizzes reinforce comprehension. Completed concepts and best quiz scores are persisted per authenticated user through the learning progress API with row-level security, with a browser-based fallback for guest users that merges safely on sign-in.

Financial education through the Learn module prepares users to understand the subsequent assessment and investment recommendation modules. A user who has completed the budgeting, emergency fund, compound interest, risk and return, mutual fund, SIP and risk tolerance concepts possesses the conceptual foundation needed to interpret the assessment outputs and the recommendation explanations.